\section{Conclusion}
\label{sec:conclusion}

In this paper, we identified and resolved the catastrophic scale mismatch inherent in non-neural retrieval for Time Series Foundation Models. By introducing a mathematical formulation that adds no trainable parameters to the retrieval path---normalizing context prior to indexing and explicitly restoring continuation magnitude post-retrieval using only the query's own statistics---we demonstrated an 85.4\% recovery in retrieval MSE and prevented a 3.8$\times$ forecasting degradation on sparse data. Search remains exact throughout, executed by a flat CPU-resident FAISS index at benchmark scale and by a bit-identical batched GPU implementation of the same exact search on the full 30,490-series M5 panel.

We are deliberate about the parameter accounting. Indexing, search, and scale restoration are entirely parameter-free, and the dense-regime configuration keeps the full pipeline so. The sparse-regime configuration adds exactly one lightweight learned component---a LightGBM gate over six origin-available diagnostic features, fitted on training origins only---which remains three orders of magnitude smaller than the neural mixers it replaces, but which we decline to describe as parameter-free.

Crucially, our empirical analysis establishes a clear \textit{Taxonomy of Retrieval Regimes}. State-of-the-art neural retrieval adapters excel on dense, continuous datasets, where their parameter overhead and retraining requirements buy real accuracy; on ETTm2 our statistical fusion is both slower and less accurate than the official neural adapter, and marginally worse than the frozen backbone it augments. For highly sparse, zero-inflated domains, or where resources preclude training a neural adapter at all, explicit mathematical scale restoration provides a robust, deployable augmentation strategy.

We also report the limits of that claim without qualification. On the held-out M5 panel, ScaleRAG improved squared-error accuracy over its frozen backbone by +5.49\% RMSSE, yet beat the strongest classical baseline by only +0.69\%---below our pre-registered 3\% threshold---lost the official WRMSSE metric to both LightGBM and Seasonal-Naive, and was surpassed by the target-only backbone on MAE, WAPE, MASE, pinball loss, and interval coverage. All three pre-registered success criteria failed. The contribution of this work is therefore a validated mechanism and a regime taxonomy: scale restoration is demonstrably necessary for statistical retrieval to function at all, while the end-to-end value of that retrieval is bounded by regime and by metric. We regard the delineation of that boundary as more useful to practitioners than a leaderboard result, and we leave calibration-preserving fusion---the clearest deficiency our evaluation exposed---to future work.
